\section{Related Work}

\textbf{Flash crowds.} Jung et al.\ characterised flash crowds and separated
them from denial-of-service attacks~\cite{jung2002flash}; Wendell and Freedman
measured four years of flash crowds in CoralCDN~\cite{wendell2011going}, and
Ari et al.\ absorbed modelled flash crowds with multi-level
caching~\cite{ari2003flash}. These
crowds are often unforeseen~\cite{wendell2011going}; a result release is a
crowd whose time is announced in advance, which changes what a portal can
prepare.

\textbf{Caching and pre-computation.} Challenger et al.\ kept cached dynamic
pages consistent with the data they were built from~\cite{challenger1999scalable}, and Labrinidis
and Roussopoulos compared serving web views virtually, materialised in the
database, and materialised at the web server~\cite{labrinidis2000webview}.
Cache stampedes and their defences, such as leases~\cite{nishtala2013memcache}
and probabilistic early recomputation (XFetch)~\cite{vattani2015xfetch}, are
well studied for keys shared by many readers.

\textbf{Overload control.} SEDA~\cite{welsh2001seda}, front-end admission
control for e-commerce sites~\cite{elnikety2004gatekeeper}, adaptive admission
control that detects flash crowds~\cite{chen2005news}, and more recent
systems such as DAGOR~\cite{zhou2018dagor} and Breakwater~\cite{cho2020breakwater}
shed or queue load inside or in front of servers. User-facing waiting rooms
that issue tickets and positions are available as products~\cite{cloudflare2026wr,aws2022vwr}.
The closest research system is WebQ~\cite{doshi2015webq}: two proxies hold
users of an overloaded server in a virtual queue, tell them how long to wait,
check with a cryptographic token that they did, and shape the load to the
server's capacity. We add a comparison with other designs, queue order, lost
tickets and refreshing students.

\textbf{Retries.} Bronson et al.\ described \emph{metastable failures}, in which
a system stays overloaded after the trigger that overloaded it is
gone~\cite{bronson2021metastable}. In a study of 22 such failures from public
incident reports, retries kept the system down in more than half of
them~\cite{huang2022metastable}, and operators are warned that retries can
destabilize a system~\cite{beyer2016sre}. We measure the effect of students
pressing refresh, not of automated retries.

\textbf{Portal load studies.} Recent studies of university academic systems
compare a plain database design against Redis cache-aside under a fixed number
of concurrent users: Putranto and Pinandito with a hosted PostgreSQL database,
Redis and up to \num{10000} JMeter users for \SI{60}{s}~\cite{putranto2026redis},
and Tofan et al.\ in a
prototype whose database and cache latencies are fixed
delays~\cite{tofan2026cacheaside}. Our study adds
pre-computed and waiting-room designs, an open-loop release-time surge with
refreshing students~\cite{schroeder2006open}, correctness checking, and
resource cost, and follows established guidance for reproducible performance
evaluation~\cite{papadopoulos2019methodological,uta2020bigdata}.
